\documentclass[11pt]{article}

\usepackage[T1]{fontenc}
\usepackage{lmodern}
\usepackage{microtype}
\usepackage{amsmath,amssymb,mathtools}
\usepackage[margin=1in]{geometry}

\setlength{\parindent}{0pt}
\setlength{\parskip}{0.8em}
\setlength{\emergencystretch}{2em}

\begin{document}

\section*{How the Proof Works}

For the global three-dimensional Navier--Stokes regularity problem, the fluid obeys
\[
\partial_tu+(u\cdot\nabla)u+\nabla p=\nu\Delta u,
\qquad
\nabla\cdot u=0,
\qquad
u(0)=u_0,
\quad \nu>0.
\]
The nonlinear term can sharpen a velocity change by carrying it into less space. Viscosity responds to that sharpening through $\Delta u$, which grows as the spatial scale shrinks, while pressure keeps the evolving field divergence-free. Concentration strengthens both the nonlinear change and the viscous response at once. If transport keeps pushing motion inward faster than viscosity can spread it back out, the velocity changes across less and less space.

A vortex makes that conflict concrete: as it stretches lengthwise, incompressibility forces its core to thin. If its velocity changes by a characteristic amount $U$ across a core of radius $r$, then
\[
|\nabla u|\sim \frac{U}{r},
\qquad
|(u\cdot\nabla)u|\sim \frac{U^2}{r},
\qquad
|\nu\Delta u|\sim \nu\frac{U}{r^2}.
\]
As the vortex stretches and thins, fluid around its inner radius can move faster and faster: $U$ may rise while $r$ falls. Holding $U$ fixed would make viscosity gain an extra power of $1/r$, but a concentrating Navier--Stokes flow does not hold $U$ fixed. The exact way to see the competition is to zoom the equation itself. For $\lambda>1$, set
\[
u_\lambda(x,t)=\lambda u(\lambda x,\lambda^2t),
\qquad
p_\lambda(x,t)=\lambda^2p(\lambda x,\lambda^2t).
\]
This zoom shrinks length by $\lambda^{-1}$, shrinks time by $\lambda^{-2}$, and raises velocity by $\lambda$. Each term in Navier--Stokes then acquires one common factor, with every expression on the right evaluated at $(\lambda x,\lambda^2t)$:
\[
\begin{aligned}
\partial_tu_\lambda&=\lambda^3\partial_tu,&
(u_\lambda\cdot\nabla)u_\lambda&=\lambda^3(u\cdot\nabla)u,\\
\nabla p_\lambda&=\lambda^3\nabla p,&
\nu\Delta u_\lambda&=\lambda^3\nu\Delta u.
\end{aligned}
\]
Nonlinear sharpening and viscous diffusion remain evenly scaled under concentration. The Sobolev measurements show what that concentration hides and what it exposes:
\[
\|u_\lambda(\cdot,t)\|_{\dot H^s}^2
=
\lambda^{2s-1}
\|u(\cdot,\lambda^2t)\|_{\dot H^s}^2.
\]
Below $s=\tfrac12$, the measured size falls as the zoom tightens. Above $s=\tfrac12$, it rises. At the critical height $s=\tfrac12$, it does neither. In particular,
\[
\|u_\lambda\|_2^2=\lambda^{-1}\|u\|_2^2,
\]
so finite kinetic energy can coexist with gradients that become larger and larger. Energy sees the total motion, but it can miss how tightly that motion is packed.

The scaling also gives the short time available at each scale. Suppose the active radius contracts by a fixed factor $0<\theta<1$ from one stage to the next,
\[
r_j=\theta^j r_0.
\]
The parabolic scaling assigns that stage a time of order $\Delta t_j\asymp r_j^2/\nu$. Hence
\[
\sum_{j=0}^{\infty}\Delta t_j
\asymp
\frac{r_0^2}{\nu}
\sum_{j=0}^{\infty}\theta^{2j}
=
\frac{r_0^2}{\nu(1-\theta^2)}
<\infty.
\]
This finite geometric sum is the Zeno cascade: infinitely many smaller and faster scale changes can fit inside a finite interval while the kinetic energy remains finite. This calculation identifies the threat. The proof has to rule it out inside the exact velocity--pressure history generated by $u_0$.

In order for a singularity to happen, velocity has to blow up. In order for that to happen, its acceleration would also have to blow up, which means its jerk has to blow up, and so on and so on:
\[
u,\qquad \partial_tu,\qquad \partial_t^2u,\qquad \partial_t^3u,\qquad\ldots
\]
What you would expect to see in this pattern, if a singularity were to form, is less spatial regularity as you climb the time derivatives, not more. Navier--Stokes reverses that expectation. The first three differentiated equations let the reversal appear before any shorthand is introduced:
\[
\begin{aligned}
\partial_tu={}&\nu\Delta u-(u\cdot\nabla)u-\nabla p,\\
\partial_t^2u={}&\nu\Delta\partial_tu
-(\partial_tu\cdot\nabla)u
-(u\cdot\nabla)\partial_tu
-\nabla\partial_tp,\\
\partial_t^3u={}&\nu\Delta\partial_t^2u
-(\partial_t^2u\cdot\nabla)u
-2(\partial_tu\cdot\nabla)\partial_tu
-(u\cdot\nabla)\partial_t^2u
-\nabla\partial_t^2p.
\end{aligned}
\]
The nonlinear terms expand by the product rule, pressure differentiates with them, and the Laplacian moves one level upward through the temporal sequence. The complete pattern is
\[
\begin{aligned}
\partial_t^{\,n+1}u={}&\nu\Delta\partial_t^{\,n}u
-\sum_{a=0}^{n}\binom{n}{a}
\bigl(\partial_t^{\,a}u\cdot\nabla\bigr)
\partial_t^{\,n-a}u
-\nabla\partial_t^{\,n}p.
\end{aligned}
\]
Writing
\[
V_n:=\partial_t^{\,n}u,
\qquad
Q_n:=\partial_t^{\,n}p
\]
now compresses a pattern the reader has already seen:
\[
V_{n+1}
+\sum_{a=0}^{n}\binom{n}{a}(V_a\cdot\nabla)V_{n-a}
+\nabla Q_n
=\nu\Delta V_n.
\]
Viscosity has one exact role here. The term $\nu\Delta V_n$ connects the next temporal response to two further spatial derivatives inside the preceding response. Transport and pressure travel with that connection as a complete coupled row; neither can be dropped to make the relation look easier.

At the critical height, that connection turns into a visible Sobolev ladder. Put
\[
E_s(t):=\frac12\|\Lambda^su(t)\|_2^2,
\qquad
H(t):=E_{1/2}(t),
\]
and keep the full signed nonlinear--pressure transfer
\[
\mathcal P_s
:=-\left\langle
\Lambda^su,
\Lambda^s\bigl((u\cdot\nabla)u+\nabla p\bigr)
\right\rangle.
\]
The energy identity is
\[
E_s'=-2\nu E_{s+1}+\mathcal P_s.
\]
Starting at $s=\tfrac12$ and differentiating gives
\[
\begin{aligned}
H={}&E_{1/2},\\
H'={}&-2\nu E_{3/2}+\mathcal P_{1/2},\\
H''={}&(-2\nu)^2E_{5/2}
+\partial_t\mathcal P_{1/2}
-2\nu\mathcal P_{3/2},
\end{aligned}
\]
and, for every $n\ge1$,
\[
H^{(n)}
=(-2\nu)^nE_{n+1/2}
+\sum_{j=0}^{n-1}(-2\nu)^j
\partial_t^{\,n-1-j}\mathcal P_{j+1/2}.
\]
So
\[
H,\ H',\ H'',\ H''',\ldots
\]
expose
\[
E_{1/2},\ E_{3/2},\ E_{5/2},\ E_{7/2},\ldots.
\]
The higher temporal responses do not point toward lower spatial regularity. Once the complete rows are controlled, they point upward through the spatial Sobolev tower. This is the reversal:
\[
\text{higher temporal order}
\longrightarrow
\text{higher spatial Sobolev order}.
\]
Its payoff is immediate. For any integer $k\ge0$,
\[
s>k+\frac32
\qquad\Longrightarrow\qquad
H^s(\mathbb R^3)\hookrightarrow C^k(\mathbb R^3).
\]
If every finite spatial height is present, then every finite $k$ is present as well:
\[
\bigcap_{s<\infty}H^s(\mathbb R^3)
=H^\infty(\mathbb R^3)
\hookrightarrow C^\infty(\mathbb R^3).
\]
The temporal sequence that looked as though it should expose worsening spatial roughness is pointing toward smoothness instead.

The remaining proof payment is to control those rows without pretending that one estimate controls the entire infinite tower. Every actual analytical demand is finite. Fix a temporal cutoff $N$ and a spatial cutoff $M$. The whole-terminal rectangle theorem gives, for every finite pair $n,m$ in that block,
\[
V_n\in L^2(0,T_*;H^{m+1}),
\qquad
V_{n+1}\in L^2(0,T_*;H^{m-1}),
\]
with a finite whole-terminal bound for that rectangle. Its constant may depend on $N$ and $M$; no uniform constant across the tower is needed. Pressure remains attached to the velocity coordinates through
\[
-\Delta Q_n
=
\partial_i\partial_j\partial_t^n(u_i u_j),
\]
and the differentiated recurrence constructs every coordinate from the one datum-generated history.

Two finite rectangles overlap on finitely many coordinates. On that overlap they agree, because each coordinate is a derivative of the velocity and pressure generated by $u_0$. Compatibility lets all finite rectangles assemble into the projective intersection
\[
u_0
\longrightarrow
\{\mathcal R_{N,M}[u_0]\}_{N,M<\infty}
\longrightarrow
\mathcal I[u_0]
:=
\bigcap_{N,M<\infty}H_t^N(0,T_*;H_x^M).
\]
The conditional Sobolev payoff is now an actual property of the solution: every finite temporal order and every finite spatial height is available inside $\mathcal I[u_0]$.

That intersection reaches the endpoint. For every finite $m$ it contains the adjacent pair
\[
u\in L^2(0,T_*;H^{m+1}),
\qquad
u_t\in L^2(0,T_*;H^{m-1}).
\]
The Lions--Magenes trace theorem then gives
\[
u\in C([0,T_*];H^m)
\qquad\text{for every finite }m,
\]
so the terminal value exists and satisfies
\[
u_*:=u(T_*)\in\bigcap_{m<\infty}H^m=H^\infty.
\]
The pressure equation and the binomial recurrence keep every temporal response in $H^\infty$ as well. If $T_*<\infty$, smooth local existence restarts the solution from $u_*$, and uniqueness glues that continuation to the history already constructed. The supposed maximal time was not maximal. Thus $T_*=\infty$.

The familiar estimates come after the intersection because they can now be selected at whatever finite height they need. Given finite $r$ and $k$, and $2\le p\le\infty$, choose
\[
m>k+\frac32-\frac3p.
\]
Sobolev embedding then gives the requested finite norm of $\nabla^k\partial_t^ru$. In particular,
\[
\mathcal I[u_0]\longrightarrow C_tH_x^1
\longrightarrow L_t^\infty L_x^6,
\qquad
\mathcal I[u_0]\longrightarrow C_tH_x^2
\longrightarrow L_t^\infty L_x^\infty,
\]
and
\[
\mathcal I[u_0]\longrightarrow C_tH_x^3
\longrightarrow L_t^\infty W_x^{1,\infty}.
\]
No one estimate carries the proof. Each finite estimate is drawn from an intersection that is already present. The whole route is
\[
\begin{aligned}
u_0
&\longrightarrow \{(V_n,Q_n)\}_{n\ge0}
\longrightarrow \{\mathcal R_{N,M}[u_0]\}_{N,M<\infty}
\longrightarrow \mathcal I[u_0]\\
&\longrightarrow u_*\in H^\infty
\longrightarrow T_*=\infty.
\end{aligned}
\]

\end{document}
